% !PRINTED 747
\weberpage{747}
In order to prove this important theorem, we recall that every algebraic number $\Theta$ is the root of one and only one irreducible equation
\weq{1}{f(\Theta) = a_0 \Theta^n + a_1 \Theta^{n-1} + \cdots + a_n = 0}
in which $a_0, a_1, \ldots, a_n$ are integral rational numbers without a common divisor, and $a_0$ is different from zero and positive. The degree $n$ of equation (1) is a positive integer, and therefore at least $=1$. The numbers $a_1, \ldots, a_{n-1}$ may in part be zero.

We now wish to determine the signs $\pm$ in such a way that $\pm a_1, \pm a_2, \pm a_3, \ldots, \pm a_n$ are not negative, and we call the sum
\weq{2}{N = (n-1) + a_0 \pm a_1 \pm a_2 \cdots \pm a_n}
the \gesperrt{height} of the algebraic number $\Theta$. The height is then always a \gesperrt{positive} integer.

Now it is easy to see that for a given value of the height $N$ there can always be only a finite number of algebraic numbers. For, in the first place, by (2) $n$ can never be greater than $N$, and for each given $N$ and $n$ the numbers $a_0, a_1, \ldots, a_n$ can be determined in only a finite number of ways. One has then to retain, among the functions $f(\Theta)$ so determined, only the irreducible ones. If we now order the algebraic numbers in such a way that we place the numbers of smaller height before those of greater height, that among numbers of equal height we place first those whose real part is smaller, and that among numbers of equal height and equal real part we let those of smaller imaginary part precede, then we have a countable arrangement of the algebraic numbers, and it is proved that the totality of all algebraic numbers is an enumerable set.

We obtain, for example:
\[
\begin{array}{llll}
N=1, & n=1, & a_0=1, & a_1=0\\
N=2, & n=1, & a_0=1, & a_1=\pm 1\\
N=3, & n=1, & a_0=1, & a_1=\pm 2\\
     &      & a_0=2, & a_1=\pm 1\\
     & n=2, & a_0=1, & a_1=0, \quad a_2=1,
\end{array}
\]
and the beginning of the ordered sequence of the algebraic numbers therefore becomes:
\[
0, -1, +1, -2, -\tfrac{1}{2}, -\sqrt{-1}, \sqrt{-1}, \tfrac{1}{2}, 2, \ldots
\]
